\chapter{Methodology}
\label{cha:methodology}

This thesis investigates how step-level procedural descriptions can be transformed into explicit procedural models that support reasoning, validation, and traceability. The input to the methodology is a set of temporally bounded step segments derived from task recordings. Each segment may include operator speech, manual annotations, and optional perception-derived evidence such as detected objects or spatial relations.

The purpose of the methodology is not to infer the complete procedure directly from raw video. Instead, the focus is on representing step-level evidence symbolically and using rule-based reasoning to make implicit procedural knowledge explicit. This includes inferred preconditions, expected effects, safety requirements, tool and material relations, and explanation traces.

The resulting representation supports the construction of a procedural graph in which the given step sequence is preserved, while additional reasoning links describe dependencies, supporting evidence, and inferred constraints.

\section{The Reasoning Pipeline}
\label{sec:method-reasoning-pipeline}

The proposed methodology is organized as a layered reasoning pipeline. Each layer performs a specific transformation, from step-level input evidence to symbolic facts, inferred constraints, validated steps, and finally a procedural graph.

The pipeline separates evidence preparation from reasoning. The early layers organize information associated with each step segment, including textual descriptions, annotations, semantic labels, and optional spatial relations. The later layers apply explicit rules to infer procedural conditions and evaluate whether each step is sufficiently supported.

This separation is important because the sources of evidence may be uncertain or incomplete, while the reasoning process itself is rule-based. Given the same symbolic facts, rules, thresholds, and ontology definitions, the reasoning layer produces the same inferred constraints, validation outcomes, and graph structure.

The four layers are:

\begin{itemize}
    \item Layer 1 defines the step-aligned input package.
    \item Layer 2 converts the input package into symbolic facts.
    \item Layer 3 applies rule-based inference to derive implicit procedural constraints.
    \item Layer 4 validates and enriches the steps and constructs the procedural graph.
\end{itemize}

The following sections define the notation, assumptions, and transformations used in each layer.


\subsection{Definitions and notation for reasoning elements}
\label{subsec:method-definition-notation-reasoning}

The methodology uses a small set of reasoning elements that are shared between layers, to clarify the relationship between the proposed architecture and prior formulations such as \citet{Zhao2025VideoToBT}, a lightweight notation is introduced to describe the elements used across layers.

\begin{itemize}
 \item $\mathcal{O}$: a set of relevant objects present or mentioned in the steps' scenes.
 \item $\mathcal{S}$: a set of candidate step segments provided as input.
 \item $\mathcal{P}$: a set of unary and binary predicates representing step-level, semantic, and spatial evidence.
 \item $\mathcal{C}$: a set of inferred constraints representing assembly-relevant conditions.
 \item $\mathcal{A}$: a set of abstract action types.
 \item $\mathcal{H}$: a set of validated or candidate procedural hypotheses.
\end{itemize}


The term candidate does not mean that the reasoning layer generates the step order. It means that each provided step description is treated as a hypothesis whose content, requirements, and effects must be validated.

An entity $o_j \in O$ may correspond to a material, part, tool, fixture, or workstation element. Object identities may originate from manual annotations, transcripts, prior domain knowledge, or perception models. When several sources provide conflicting labels, they are represented as competing semantic hypotheses.

Predicates represent symbolic evidence. They may describe step-level information, semantic identity, or spatial relations. Examples include $hasAction(s_i,place)$, $usesObject(s_i,o_j)$, $usesTool(s_i,o_k)$, $isA(o_j,Rod)$, and $aligned(o_j,o_k)$. Each predicate includes a confidence value and provenance information linking it to its source.

Constraints represent procedural or assembly-relevant conditions inferred from predicates. These include preconditions, expected effects, safety requirements, tool requirements, compatibility relations, and implicit assembly conditions. Constraints are not directly observed. They are produced by rules and are linked to the predicates that triggered them.

Rules define how predicates are mapped into constraints. Inference rules produce supporting constraints, while compatibility rules identify action-object combinations that are inadmissible. This distinction allows the system to separate positive evidence from hard violations.

For each step $s_i \in \mathcal{S}$ associated with a time interval in the recording, the system operates on a predicate set $\mathcal{P}_{s_i}$ and an inferred constraint set $\mathcal{C}_{s_i}$ associated with source descriptions and optional evidence. Layer 4 uses these sets to validate and enrich the step and the final procedural graph is denoted by

\[
G = (V,E),
\]

where $V$ contains validated step nodes and associated evidence nodes, and $E$ contains temporal, dependency, and support relations.



\section{Methodological overview and module boundaries}
\label{sec:method-overview}

The system is divided into a front-end and a backend. The front-end prepares symbolic evidence from step-aligned inputs, while the backend performs rule-based inference, validation, enrichment, and graph construction.

\begin{itemize}
    \item Layer 1: Step-aligned input package
    \begin{itemize}
        \item Defines candidate step segments and temporal boundaries.
        \item Associates each step with transcripts, annotations, and optional perception evidence.
        \item Preserves source metadata, confidence values, and temporal references.
    \end{itemize}

    \item Layer 2: Symbolic fact construction
    \begin{itemize}
        \item Converts step descriptions and available evidence into symbolic predicates.
        \item Represents actions, objects, tools, materials, quantities, and time windows.
        \item Adds semantic and spatial predicates when such evidence is available.
        \item Maintains confidence, provenance, and temporal references for generated facts.
    \end{itemize}

    \item Layer 3: Rule-based procedural inference
    \begin{itemize}
        \item Evaluates domain rules over the symbolic fact base.
        \item Infers preconditions, expected effects, safety requirements, tool requirements, compatibility constraints, and implicit assembly conditions.
        \item Propagates confidence values through rule evaluation.
        \item Links inferred constraints to the predicates and rules that produced them.
    \end{itemize}

    \item Layer 4: Step validation and procedural graph construction
    \begin{itemize}
        \item Evaluates each step against predicates and inferred constraints.
        \item Classifies step support as accepted, uncertain, or rejected.
        \item Enriches step nodes with inferred procedural information.
        \item Constructs a procedural graph containing temporal, dependency, and support relations.
    \end{itemize}
\end{itemize}

Layers 1 and 2 define the front-end interface. They organize heterogeneous sources into a symbolic representation without depending on a specific perception model. Layers 3 and 4 define the backend contribution of the thesis, where explicit rules are used to infer, validate, and structure procedural knowledge.

The output is a procedural graph that preserves the provided step sequence while adding reasoning-derived relations. The graph is therefore not only a transcription of the input descriptions, but a structured representation of explicit and inferred procedural knowledge.


\subsection{Input assumptions and domain preparation}
\label{sec:input-assumptions}

This section defines the assumptions and conventions used before the reasoning layers are applied. The purpose is to ensure that step records, entities, predicates, and rules share a common structure.

The input procedure is represented as an ordered set of step segments $\mathcal{S} = \{s_1,\ldots,s_n\}$. Each step $s_i$ has a time interval, one or more source descriptions, and optional supporting evidence. The order of the step segments is assumed to be given by the input annotations or segmentation process.

The pipeline operates over a domain-configurable set of entity types, such as rods, sleeves, supports, tools, holes, or workstations. These entity types are not fixed by the methodology. They are defined in the domain configuration and can be adapted to different assembly scenarios.

Each relevant entity mentioned or observed during a step is assigned an identifier. The same entity may appear across multiple steps when there is sufficient evidence to link it. For example, a rod mentioned in one step may later appear as the object on which sleeves are placed. These links may come from annotations, transcripts, perception outputs, or prior domain knowledge.

The predicate vocabulary is kept compact and explicit. Step-level predicates describe actions and involved entities, for example $hasAction(s_i,place)$ and $usesObject(s_i,o_j)$. Semantic predicates describe entity types, for example $isA(o_j,Rod)$. Spatial predicates describe qualitative relations between entities, for example $aligned(o_j,o_k)$ or $touching(o_j,o_k)$, when geometric evidence is available.

The domain configuration also defines the rule vocabulary used by Layer 3. Rules may encode action compatibility, required preconditions, expected effects, tool requirements, safety requirements, and implicit assembly conditions. This allows the reasoning layer to infer information that is often missing from speech or manual annotations.

All intermediate representations include provenance information. This records whether a predicate or constraint originated from operator speech, manual annotation, perception output, prior domain knowledge, or rule-based inference. This provenance is necessary for explanation traces and for distinguishing explicitly stated information from inferred information.


\subsubsection{Rule set authoring and constraint semantics}
\label{sec:rule-authoring}

Given the definitions introduced in \cref{subsec:method-definition-notation-reasoning}, this section describes how rules are specified and how their outputs are interpreted within the reasoning process.
Inference rules are defined as conjunctions of qualitative predicates and semantic conditions, independent of low-level geometry. Their role is to map perceptual evidence into higher-level assembly conditions expressed as constraints. These constraints are qualitative in nature and are interpreted operationally within the reasoning context rather than as fixed numeric thresholds.

When multiple inference rules are applicable, their outputs may coexist and are associated with confidence values derived from supporting predicates. This allows the system to maintain multiple plausible interpretations of the same scene without committing to a single deterministic outcome.

Compatibility rules are defined separately and operate over semantic predicates to determine whether a given action is admissible for a pair of object types. Unlike inference rules, they do not contribute supporting constraints, but instead produce incompatibility relations that are enforced during procedural validation.

In this work, the term \textit{hypothesis} refers to candidate interpretations at different levels of abstraction, including semantic labels, spatial relations and inferred constraints. These hypotheses are not treated as ground truth, but as uncertainty-bearing representations that may coexist and be evaluated or refined across layers.


\section{Layer 1 methodology: step-aligned input package}
\label{sec:layer1-method}

Layer 1 defines the input interface between task recordings and the reasoning pipeline. The input is an ordered set of candidate step segments $\mathcal{S} = \{s_1,\ldots,s_n\}$, where each $s_i$ corresponds to a temporally bounded part of the recorded procedure.

Each step segment may include operator speech, manual annotations, and optional perception-derived evidence. Operator speech provides the description given during task execution, while manual annotations may correct, refine, or add details after the recording has been reviewed. Optional perception evidence may include detected entities, semantic labels, poses, or geometric representations.

The output of Layer 1 is a structured record for each step. A record includes the step identifier, sequence index, time interval, source descriptions, and any available entity or perception references. This representation preserves the link between the symbolic reasoning process and the original recording.

Layer 1 does not validate the correctness or completeness of the step description. It only organizes the available evidence into a consistent input structure. Validation, enrichment, and inference are performed in later layers.


\section{Layer 2 methodology: symbolic fact construction}
\label{sec:layer2-method}

Layer 2 transforms each step-aligned input record into a symbolic fact base. For each step $s_i$, the output is a local predicate set $\mathcal{P}_{s_i}$ containing step-level, semantic, and, when available, spatial predicates.

Step-level predicates describe the procedural information associated with the step. These include the action, time window, involved objects, tools, materials, and quantities. For example, a step in which the operator places copper sleeves on a rod may produce predicates such as $hasAction(s_i,place)$, $usesObject(s_i,o_1)$, $usesObject(s_i,o_2)$, $hasQuantity(o_1,8)$, and $hasTimeWindow(s_i,t_a,t_b)$.

Semantic predicates describe the identity or role of entities. These may be derived from manual annotations, transcripts, prior domain knowledge, or perception outputs. For example, Layer 2 may produce $isA(o_1,Sleeve)$, $isA(o_2,Rod)$, or $isA(o_3,Screwdriver)$. When multiple interpretations are possible, competing semantic hypotheses may be retained with separate confidence values.

When geometric evidence is available, Layer 2 may also derive qualitative spatial predicates. Oriented bounding boxes can be mapped to relations such as \textit{aligned}, \textit{inside}, and \textit{touching}, following a compact spatial abstraction based on oriented bounding boxes rather than full meshes \citep{Haesler_2025}. These predicates are optional evidence for reasoning and are not required for every step.

All predicates include confidence and provenance information. Provenance records whether a predicate originated from operator speech, manual annotation, perception output, prior domain knowledge, or another preprocessing source. This allows later layers to distinguish explicitly stated facts from uncertain or automatically derived evidence.

As illustrated later in \cref{sec:example-rule1}, Layer 2 may produce both step-level predicates and spatial predicates for the same step. For example, a candidate insertion step may include $hasAction(s_i,insert)$, $usesObject(s_i,o_1)$, $usesObject(s_i,o_2)$, and $isA(o_1,Peg)$, together with optional spatial evidence such as $aligned(o_1,o_2)$ and $inside(o_1,o_2)$. These predicates form the traceable input for Layer 3.

If upstream geometry is unstable, Layer 2 could optionally be preceded by a representation or reconstruction stage that improves local geometric consistency before relation extraction, for example by reconstructing scenes from primitive-level elements \citep{Mazur_2024}. In the present thesis, this is treated as an optional extension rather than a core contribution.


\section{Layer 3 methodology: rule-based procedural inference}
\label{sec:layer3-method}

Layer 3 performs rule-based inference over the predicate sets produced in Layer 2. For each step $s_i$, the layer consumes the local predicate set $\mathcal{P}_{s_i}$ and derives a set of inferred constraints $\mathcal{C}_{s_i}$. These constraints make implicit procedural knowledge explicit, including preconditions, expected effects, safety requirements, required tools, compatibility conditions, and implicit assembly conditions.

\subsection{Conceptual design}

The input to Layer 3 consists of symbolic predicates associated with a step. These predicates may describe the action performed in the step, the involved objects, their semantic types, and optional spatial relations. Layer 3 interprets these predicates using explicit domain rules.

A rule is defined as a conjunction of unary and binary predicates. When the antecedents of a rule are satisfied, the system infers one or more higher-level constraints. Each inferred constraint is associated with a confidence value derived from the supporting predicates and with provenance information linking it to the rule and evidence that produced it.

For a step $s_i$, let $\mathcal{P}_{s_i}$ denote the local predicate set available for reasoning. A rule $r$ maps a subset of supporting predicates $\mathbf{p}_r \subseteq \mathcal{P}_{s_i}$ to one or more inferred constraints:

\[
r : \mathbf{p}_r \rightarrow \mathbf{c}_r,
\qquad
\mathbf{c}_r \subseteq \mathcal{C}_{s_i}.
\]

The rule is activated when all supporting predicates satisfy the rule threshold:

\[
\min_{p \in \mathbf{p}_r} \text{conf}(p) \geq \tau.
\]

The confidence of each inferred constraint is computed using the same conservative aggregation policy:

\[
\text{conf}(c) = \min_{p \in \mathbf{p}_r} \text{conf}(p).
\]

This formulation keeps the inference mechanism uniform across different rule types. The difference between rule types lies in the semantic role of the inferred constraint, not in the evaluation mechanism.

\paragraph{Rule types.}
Layer 3 uses several types of rules. Inferred precondition rules derive conditions that should hold before a step can be considered valid. For example, placing sleeves on a rod requires the rod to have already been placed on its supports:

\begin{align*}
&isA(o_1,\text{Sleeve}) \wedge isA(o_2,\text{Rod}) \wedge isA(o_3,\text{Support}) \\
&\wedge hasAction(s_i,\text{place}) \wedge usesObject(s_i,o_1) \wedge usesObject(s_i,o_2) \\
&\Rightarrow requires(s_i, placed(o_2,o_3)) 
\end{align*}

Safety constraint rules derive safety-related conditions that may be omitted from the step description. For example, tightening components on a rod requires the rod to be secured on the workstation:

\begin{align*}
&isA(o_1,\text{Rod}) \wedge isA(o_2,\text{Workstation}) \\
&\wedge hasAction(s_i,\text{tighten}) \wedge usesObject(s_i,o_1) \\
&\Rightarrow requiresSafety(s_i, secured(o_1,o_2))
\end{align*}

Required tool rules infer tools needed for a step. For example, tightening a sleeve requires a screwdriver or equivalent tool:

\begin{align*}
&isA(o_1,\text{Sleeve}) \wedge hasAction(s_i,\text{tighten})  \\
&\wedge usesObject(s_i,o_1)  \\
&\Rightarrow requiresTool(s_i,\text{Screwdriver})
\end{align*}

Implicit assembly condition rules derive technical conditions required for an action to make sense. For example, a sleeve can only be tightened to a rod if their corresponding holes are aligned:

\begin{align*}
&isA(o_1,\text{Sleeve}) \wedge isA(o_2,\text{Rod})  \\
&\wedge hasAction(s_i,\text{tighten}) \wedge usesObject(s_i,o_1) \wedge usesObject(s_i,o_2) \\
&\Rightarrow requires(s_i, alignedHoles(o_1,o_2))
\end{align*}

Compatibility rules identify inadmissible action-object combinations. Unlike the previous rule types, they act as hard validity conditions during step validation. For example, if an action is not meaningful for a given pair of object types, the rule produces an incompatibility relation:

\begin{align*}
&isA(o_1,T_1) \wedge isA(o_2,T_2) \wedge hasAction(s_i,a) \\
&\Rightarrow incompatibleAction(s_i,o_1,o_2,a).
\end{align*}

The resulting constraints are passed to Layer 4, where they are used to validate and enrich the corresponding step.

\Cref{fig:layer3_diagram-1} illustrates the rule-based inference process, where spatial and semantic predicates are combined with compatibility rules to produce assembly constraints and action admissibility conditions, with uncertainty propagated through confidence aggregation.

\begin{figure}[t]
  \begin{center}
    %\vspace*{1em}
    \includegraphics[width=0.9\linewidth]{Figures/Layer_3_with3objects.png}
    %\vspace*{1em}
  \end{center}
  \caption{Rule-based assembly inference in Layer 3. Spatial and semantic predicates are combined with
semantic compatibility rules to infer higher-level assembly constraints and action admissibility
conditions, with uncertainty propagated through confidence-based aggregation.}
  \label{fig:layer3_diagram-1}
\end{figure}


\subsection{Design decisions and alternatives}
\label{sub:layer3-method-design-decisions}

\subsubsection{Rule-based inference versus learned inference.}
An alternative would be to train a model to infer procedural constraints directly from step descriptions, annotations, or symbolic inputs. Such an approach could reduce the need for manually authored rules, but it would require task-specific training data and would make the reasoning process harder to inspect.

This thesis adopts a rule-based approach because the objective is not only to infer constraints, but also to explain where they came from. Rules make domain assumptions explicit, allow inferred constraints to be traced back to supporting predicates, and provide deterministic behavior once the symbolic inputs and configuration are fixed. The trade-off is that the quality of the reasoning depends on the coverage and correctness of the rule set.

\subsubsection{Local step reasoning versus global procedure inference.}
Another alternative would be to infer the complete procedural structure globally, including the order of steps, alternative paths, and long-range dependencies. This would allow richer process modeling, but would also require stronger assumptions about observation quality and procedure completeness.

In this work, the order of step segments is provided as input. Reasoning is performed locally for each step, and graph-level dependencies are added when inferred effects support later preconditions. This keeps the method tractable and aligned with the available evidence. The trade-off is that the system does not infer a complete executable process model.

\subsubsection{Constraint layer versus direct predicate validation.}
A simpler design would validate steps directly against predicates from Layer 2. However, many relevant procedural conditions are not directly stated or observed. For example, the requirement that sleeve holes must be aligned before tightening is not merely a detected spatial relation, but an assembly condition inferred from the action and involved objects.

For this reason, Layer 3 introduces an explicit constraint layer. Predicates provide evidence, while constraints represent inferred procedural knowledge. This separation makes it possible to distinguish what was stated, observed, or inferred.

\subsubsection{Confidence aggregation.}
A fully probabilistic approach could model dependencies between predicates and constraints more accurately. However, such models require assumptions and data that are outside the scope of this thesis, and they may reduce interpretability.

This thesis represents uncertainty using scalar confidence values attached to predicates and constraints. Rule confidence is computed using minimum aggregation over the supporting predicates. This conservative policy is simple, reproducible, and easy to trace, although it does not capture complex dependencies between evidence sources.

\subsubsection{Object-level versus site-level reasoning.}
Some assembly conditions may require reasoning about specific sites, such as holes, edges, or contact regions. Site-level reasoning would allow more precise rules, for example when checking whether a sleeve hole is aligned with a rod hole. However, it requires more detailed geometric information and more complex annotations.

The main representation in this thesis is object-level. Site-level predicates may be introduced when the required information is available, but they are treated as an extension rather than a core assumption.


\subsection{Implementation scope and delimitation}

The conceptual design supports several types of procedural inference, including preconditions, expected effects, safety requirements, tool requirements, compatibility constraints, and implicit assembly conditions. The implemented prototype focuses on a limited subset of these rule types in order to evaluate the core reasoning mechanism.

Specifically, the implementation evaluates predefined rules over the predicate set $\mathcal{P}_{s_i}$ for each step $s_i$, generates inferred constraints in $\mathcal{C}_{s_i}$, propagates confidence values using a fixed aggregation policy, and records provenance links between inferred constraints and supporting predicates.

The implementation does not attempt to learn rules automatically, infer step boundaries, perform full natural language understanding, or maintain a global hypothesis structure over the whole procedure. It also does not implement detailed site-level reasoning unless the required predicates are explicitly available.

This delimitation keeps the implementation focused on the main contribution of the thesis: transforming symbolic step evidence into traceable procedural constraints that can support validation and graph construction.

\subsection{Relation to assembly knowledge representation}

Layer 3 follows the idea that assembly knowledge can be represented hierarchically, where lower-level relations support the inference of higher-level procedural conditions. In this thesis, predicates represent explicit or observed evidence, while constraints represent inferred procedural knowledge.

This distinction is important because many relevant assembly conditions are not directly stated in the source descriptions. For example, the requirement that holes must be aligned before tightening is inferred from the action and involved objects, rather than simply copied from the annotation. This is consistent with assembly knowledge representations in which process constraints are derived from structural and semantic relations \citep{Zhu_2022}.


\section{Layer 4 methodology: step validation and procedural graph construction}
\label{sec:layer4-method}

Layer 4 validates and enriches the step segments using the predicates from Layer 2 and the constraints inferred in Layer 3. For each step $s_i$, the layer consumes the predicate set $\mathcal{P}_{s_i}$ and the constraint set $\mathcal{C}_{s_i}$, and produces a validated step representation.

\subsection{Conceptual design}

Layer 4 evaluates whether the information associated with a step is sufficiently supported by the available evidence and inferred constraints. Although the step sequence is provided as input, the content of each step is not assumed to be complete or fully correct. A spoken or manually annotated step may omit preconditions, safety requirements, required tools, or expected effects.

For each step $s_i$, the validation process considers:

\begin{itemize}
    \item the action described in the step,
    \item the objects, materials, and tools associated with the step,
    \item inferred preconditions and expected effects,
    \item safety and compatibility constraints,
    \item missing or conflicting evidence.
\end{itemize}

Let $S_p \subseteq \mathcal{P}_{s_i}$ denote the predicates relevant to the step, and let $S_c \subseteq \mathcal{C}_{s_i}$ denote the inferred constraints relevant to the step. The aggregated confidence of the validation evidence is computed as:

\[
\text{conf}(s_i) = \min_{x \in S_p \cup S_c} \text{conf}(x).
\]

A step is classified as accepted, uncertain, or rejected. A step is accepted when its action and required constraints are sufficiently supported and no incompatibility relation is present. A step is uncertain when partial support exists but some information is missing or below the acceptance threshold. A step is rejected when a hard incompatibility is present or when required support is absent.

This classification is not intended to replace the input annotation. Instead, it records the degree to which the step is supported by the available symbolic evidence and domain rules.



\begin{figure}[t]
  \begin{center}
    %\vspace*{1em}
    \includegraphics[width=0.9\linewidth]{Figures/Layer_4_simple_admisability_check.png}
    %\vspace*{1em}
  \end{center}
  \caption{Layer 4 procedural validation. Candidate step hypotheses are evaluated against observed
predicates and constraints inferred in Layer 3, including incompatibility relations derived from semantic compatibility rules,
producing explainable outcomes under uncertainty.}
  \label{fig:layer4_diagram-1}
\end{figure}




\Cref{fig:layer4_diagram-1} illustrates this decision process, where candidate step hypotheses are evaluated against predicates and inferred constraints and classified according to the decision policy defined above.


Each decision is accompanied by an explanation trace that explicitly links the outcome to the supporting predicates, inferred constraints, and any missing or conflicting evidence. This ensures that the reasoning process remains transparent and that the influence of uncertainty can be determined.

\Cref{fig:layer4_diagram-2} illustrates the construction of the procedural graph, where accepted and uncertain step hypotheses are represented as nodes. Temporal edges preserve the input step order, while additional reasoning-derived edges represent dependency, support, and provenance relations.



\begin{figure}
  \begin{center}
    %\vspace*{1em}
    \includegraphics[width=0.9\linewidth]{Figures/Layer_4_procedural_graph_construction.png}
    %\vspace*{1em}
  \end{center}
  \caption{Procedural graph construction from validated step hypotheses. Accepted and uncertain hypotheses are represented as step nodes, while rejected hypotheses may be excluded from the main procedural flow or retained for diagnostic inspection. Temporal edges preserve the input order, and additional edges represent dependencies, evidence, constraints, and provenance.}
  \label{fig:layer4_diagram-2}
\end{figure}



\subsection{Step enrichment}

In addition to validation, Layer 4 enriches each step with inferred procedural information. This includes preconditions, expected effects, safety requirements, required tools, and implicit assembly conditions produced by Layer 3.

For example, a step describing the placement of copper sleeves on a rod may be enriched with the precondition that the rod must already be positioned on its supports. A later tightening step may be enriched with the condition that the sleeves must already be placed and that the sleeve holes must be aligned with the holes of the rod.

The enriched step representation contains:

\begin{itemize}
    \item the original step identifier and time window,
    \item the action and involved entities,
    \item explicit information from the input sources,
    \item inferred preconditions and effects,
    \item inferred safety and tool requirements,
    \item validation status and confidence,
    \item an explanation trace.
\end{itemize}

\subsection{Procedural graph construction}

After validation and enrichment, Layer 4 constructs a procedural graph $G = (V,E)$. Step nodes are created from the validated step representations. The temporal order of these nodes follows the input sequence. If $s_i$ directly precedes $s_{i+1}$, the graph includes a temporal edge:

\[
NEXT(s_i,s_{i+1}).
\]

The graph may also include dependency edges. A dependency edge is added when an inferred or explicit effect of one step supports a precondition of a later step. For example, if one step produces $placed(o_1,o_2)$ and a later step requires $placed(o_1,o_2)$, then the graph records this relation as a dependency between the corresponding step nodes.

Additional nodes may represent objects, tools, preconditions, effects, safety requirements, constraints, and sources. These nodes are connected to step nodes through support and provenance relations. This makes the graph useful not only as a sequence of steps, but also as a traceable representation of how each step is supported and enriched.

The graph does not define an executable controller or a real-time guidance policy. It is an intermediate procedural representation intended for inspection, validation, and later extension.



%%%

\begin{figure}
  \begin{center}
    %\vspace*{1em}
    \includegraphics[width=0.9\linewidth]{Figures/graph_focus_on_step_1_light_mode.png}
    %\vspace*{1em}
  \end{center}
  \caption{Extract of the procedural graph constructed from validated step records. The \texttt{NEXT} edges preserve the input sequence, while \texttt{DEPENDS\_ON} edges make explicit how requirements are satisfied by effects produced by earlier steps. Node status and attached evidence indicate whether each step is accepted, uncertain, or rejected.}
  \label{fig:procedural_graph_focus_on_Step_1}
\end{figure}




%%%

\subsection{Design decisions and alternatives}
\label{sub:layer4-method-design-decisions}

\subsubsection{Generation of candidate procedure steps}

A key design decision concerns the role of Layer 4 with respect to step generation. One alternative would be to infer procedural steps directly within this layer by analyzing symbolic state transitions. While conceptually appealing, this approach introduces additional ambiguity and complexity, particularly under noisy or incomplete observations. In this work, the generation of candidate step hypotheses is treated externally as described in \cref{sec:layer2-method}, while Layer 4 focuses exclusively on their validation. This separation simplifies the reasoning process, enables controlled evaluation, and allows the layer to concentrate on consistency, uncertainty handling, and explainability.

\subsubsection{Representation of procedural knowledge}
Alternatives such as behavior trees or finite state machines provide structured representations that integrate planning and execution, as seen in frameworks such as \citep{Zhao2025VideoToBT}. However, these approaches introduce control-oriented semantics that are not required in the present context. In this work, each step is represented in terms of its action, arguments, inferred preconditions, and expected effects. This representation focuses on the semantic interpretation of the step rather than its execution, and aligns naturally with the predicate-based reasoning established in Layers 2 and 3. The trade-off is that execution strategies and control flow are not explicitly modeled.

\subsubsection{Handling of uncertainty}

The reasoning process uses confidence values to avoid reducing uncertain evidence to purely binary decisions. At the same time, it does not implement a fully probabilistic model, since that would require explicit modeling of dependencies between predicates, constraints, and step hypotheses.

Instead, uncertainty is handled locally at the step level. For each provided step segment, the system may evaluate one or more candidate interpretations supported by different subsets of predicates and constraints. These interpretations are assessed independently using the available evidence for that step and the effects produced by earlier non-rejected steps. The system does not maintain global hypothesis trajectories over time, as in inference-tree approaches such as \citet{WangEtAl2018}.

Confidence values are propagated from supporting predicates to inferred constraints and validation outcomes. They are then used to classify each step hypothesis as accepted, uncertain, or rejected. This design preserves ambiguity where evidence is incomplete while keeping the reasoning process interpretable and computationally tractable.


\subsubsection{Consistency-checking mechanism independent of the source}
A further decision concerns the independence of the validation process from the source of candidate steps. Whether steps are produced by automated perception models or by manual annotation, they are treated uniformly as hypotheses that require validation. This ensures that the reasoning process is robust not only to model errors, but also to incomplete or inconsistent human descriptions.

\subsubsection{Decoupling between perception and reasoning}
In tightly integrated systems such as \citep{Zhao2025VideoToBT}, planning, perception, and execution are combined in a closed loop. In contrast, the present work adopts a layered architecture in which Layer 4 operates exclusively on symbolic inputs from Layers 2 and 3. This separation improves modularity, allows independent evaluation of reasoning components, and facilitates controlled experimentation using synthetic inputs. The trade-off is that real-time feedback and dynamic re-planning are not addressed within this layer.

\subsubsection{Procedural graph construction vs execution-oriented representations}

A design decision concerns how validated procedural knowledge should be represented after step-level reasoning.

One alternative is to adopt execution-oriented representations such as behavior trees or finite state machines, as used in frameworks such as \citep{Zhao2025VideoToBT}. These representations provide explicit control flow, reactivity, and support for execution in robotic systems. However, they also introduce additional semantics related to scheduling, control, and recovery that are not required for the purposes of this work. Moreover, they typically assume a well-defined and consistent procedural structure, which may not hold under uncertain and partially observed conditions.

In this work, a simplified procedural graph is adopted, in which nodes correspond to validated step hypotheses and edges capture relations added during validation, such as temporal adjacency, dependency support, and provenance. This representation is derived after step validation and is not intended to function as a complete executable process model. Instead, it provides an inspection and explanation layer that can support operator guidance by exposing the validation status of each step, the evidence on which this status is based, the dependencies on previously produced effects, and any missing or incompatible conditions identified during reasoning.

The choice of a graph-based representation allows the system to preserve the given step order while representing validation status, uncertainty, dependencies, and alternative interpretations of step content. At the same time, it remains compatible with the local hypothesis evaluation approach adopted in this work, as it does not require maintaining or updating global hypothesis structures over time.

The trade-off is that the resulting representation cannot be directly executed or used for control without further transformation. However, this limitation is consistent with the scope of the thesis, which focuses on reasoning, validation, and explainability rather than execution. The procedural graph may therefore be interpreted as an intermediate artifact that could, in future work, be compiled into execution-oriented representations such as behavior trees if required.

\subsection{Implementation scope and delimitations}

The implementation focuses on local validation and graph construction. It assumes that step segments and their temporal order are provided as input. It evaluates predefined rules, validates each step against relevant predicates and constraints, and exports a graph containing temporal, dependency, support, and provenance relations.

The implementation does not perform automatic step segmentation, full action recognition, natural language understanding, or execution monitoring. It also does not revise the global step sequence. These aspects are outside the scope of the thesis and are left for future work.

\subsection{Relation to assembly reasoning}

In this framework, procedural reasoning is formulated as the validation of candidate step hypotheses against an evolving symbolic representation of the assembly. Layer 2 provides a perception-grounded description of objects and their relations, while Layer 3 infers assembly-relevant constraints from these observations. Layer 4 then evaluates whether the provided step descriptions are consistent with the available symbolic evidence, inferred constraints, and effects produced by earlier non-rejected steps.

This perspective differs from approaches that rely on predefined procedural templates or explicit state machines \citep{Stanescu_2022,Stanescu_2023,Stanescu_2024}. Instead, procedural knowledge is treated as a set of validated hypotheses grounded in observed interactions and inferred constraints. As a result, the system can reason about assembly processes even when the underlying procedure is uncertain or partially specified, while maintaining explicit traceability of its decisions.


\section{Example Rule 1: sleeve placement and tightening}
\label{sec:example-rule1}

This example illustrates how the proposed methodology operates over step-aligned inputs. The example is based on a task in which a rod is placed on supports, copper sleeves are positioned on the rod, and the sleeves are then tightened.

\subsection{Layer 1: step-aligned input}

Assume that the input contains three step segments:

\begin{itemize}
    \item $s_1$: time window $[t_1,t_2]$, description: ''First I place the rod on the supports.''
    \item $s_2$: time window $[t_2,t_3]$, description: ''Then I take eight copper sleeves and place them on the respective sections of the rod.''
    \item $s_3$: time window $[t_3,t_4]$, description: ''I tighten each sleeve with a screwdriver.''
\end{itemize}

These descriptions may come from operator speech, manual annotations, or a combination of both. At this stage, they are treated as structured input evidence, not as fully validated procedural knowledge.

\subsection{Layer 2: symbolic fact construction}

Layer 2 converts each step segment into symbolic predicates. For the first step, the following predicates may be produced:

\[
hasAction(s_1,\text{place}) \wedge usesObject(s_1,o_1) \wedge usesObject(s_1,o_2)
\]

\[
isA(o_1,\text{Rod}) \wedge isA(o_2,\text{Support})
\]

For the second step, Layer 2 may produce:

\[
hasAction(s_2,\text{place}) \wedge usesObject(s_2,o_3) \wedge usesObject(s_2,o_1)
\]

\[
isA(o_3,\text{Sleeve}) \wedge isA(o_1,\text{Rod}) \wedge hasQuantity(o_3,8)
\]

For the third step, Layer 2 may produce:

\[
hasAction(s_3,\text{tighten}) \wedge usesObject(s_3,o_3) \wedge usesTool(s_3,o_4)
\]

\[
isA(o_3,\text{Sleeve}) \wedge isA(o_4,\text{Screwdriver})
\]

If geometric evidence is available, additional spatial predicates may be added, for example:

\[
aligned(o_3,o_1), \qquad touching(o_3,o_1).
\]

These predicates form the local fact sets $\mathcal{P}_{s_1}$, $\mathcal{P}_{s_2}$, and $\mathcal{P}_{s_3}$.

\subsection{Layer 3: rule-based inference}

Layer 3 evaluates rules over the predicates associated with each step. For $s_1$, a rule may infer the expected effect of placing the rod on the supports:

\[
\begin{aligned}
&hasAction(s_1,\text{place}) \wedge isA(o_1,\text{Rod}) \\
&\wedge isA(o_2,\text{Support}) \\
&\Rightarrow produces(s_1, placed(o_1,o_2)).
\end{aligned}
\]

For $s_2$, a precondition rule may infer that sleeves can be placed only after the rod has been positioned on the supports:

\[
\begin{aligned}
&hasAction(s_2,\text{place}) \wedge isA(o_3,\text{Sleeve}) \\
&\wedge isA(o_1,\text{Rod}) \\
&\Rightarrow requires(s_2, placed(o_1,o_2)).
\end{aligned}
\]


The same step may also produce the effect that the sleeves are placed on the rod:

\[
\begin{aligned}
&hasAction(s_2,\text{place}) \wedge isA(o_3,\text{Sleeve}) \\
&\wedge isA(o_1,\text{Rod}) \\
&\Rightarrow produces(s_2, placed(o_3,o_1)).
\end{aligned}
\]


For $s_3$, Layer 3 may infer that tightening requires the sleeves to be already placed:
\[
\begin{aligned}
&hasAction(s_3,\text{tighten}) \wedge isA(o_3,\text{Sleeve}) \\
&\Rightarrow requires(s_3, placed(o_3,o_1)).
\end{aligned}
\]

It may also infer an implicit assembly condition:
\[
\begin{aligned}
&hasAction(s_3,\text{tighten}) \wedge isA(o_3,\text{Sleeve}) \\
&\wedge isA(o_1,\text{Rod}) \\
&\Rightarrow requires(s_3, alignedHoles(o_3,o_1)).
\end{aligned}
\]

Finally, a safety rule may infer that the rod should be secured before tightening:
\[
\begin{aligned}
&hasAction(s_3,\text{tighten}) \wedge isA(o_1,\text{Rod}) \\
&\Rightarrow requiresSafety(s_3, secured(o_1,\text{Workstation})).
\end{aligned}
\]


Each inferred constraint is assigned a confidence value derived from its supporting predicates and is linked to the rule that produced it.

\subsection{Layer 4: validation and enrichment}

Layer 4 validates each step using the predicates from Layer 2 and the constraints inferred in Layer 3.

For $s_1$, the action and objects are explicitly supported by the input description. The step is enriched with the expected effect $placed(o_1,o_2)$, meaning that the rod is positioned on the supports.

For $s_2$, the inferred precondition $placed(o_1,o_2)$ can be linked to the effect of $s_1$. The step is therefore supported by the previous step and enriched with the effect $placed(o_3,o_1)$, meaning that the sleeves are placed on the rod.

For $s_3$, the inferred precondition $placed(o_3,o_1)$ can be linked to the effect of $s_2$. The step is also enriched with two inferred requirements. The first is the implicit assembly condition:
\[
alignedHoles(o_3,o_1).
\]
The second is the safety requirement:
\[
secured(o_1,\text{Workstation}).
\]
If these conditions are not explicitly stated in the source descriptions, they remain inferred constraints and may be marked as uncertain or requiring confirmation.

\subsection{Graph representation and explanation trace}

The resulting procedural graph preserves the input sequence:

\[
NEXT(s_1,s_2), \qquad NEXT(s_2,s_3).
\]

In addition, reasoning-derived dependency edges are added. Since $s_2$ requires an effect produced by $s_1$, the graph contains:

\[
DEPENDS\_ON(s_2,s_1).
\]

Similarly, since $s_3$ requires the sleeves to be placed on the rod, and this condition is produced by $s_2$, the graph contains:

\[
DEPENDS\_ON(s_3,s_2).
\]

The graph also connects each step to its objects, tools, preconditions, effects, safety requirements, and supporting sources. For example, $s_3$ is connected to the screwdriver as a tool, to the sleeves as the manipulated object, to $alignedHoles(o_3,o_1)$ as an implicit assembly condition, and to $secured(o_1,\text{Workstation})$ as a safety requirement.

This example shows how the pipeline preserves the provided step order while adding reasoning-derived structure. The final graph is therefore more than a sequence of annotated steps: it records what each step requires, what it produces, and which parts of the representation were explicit, inferred, or uncertain.


\section{Example Rule 2: connector placement and securing}
\label{sec:example-rule2}

This second example illustrates how the same reasoning mechanism can be applied to a different assembly context. Consider a transformer assembly task in which a connector is placed on a terminal and then secured.

\subsection{Layer 1: step-aligned input}

Assume that the input contains two step segments:

\begin{itemize}
    \item $s_4$: time window $[t_4,t_5]$, description: ''I place the connector on the terminal.''
    \item $s_5$: time window $[t_5,t_6]$, description: ''I secure the connector with the fastening screw.''
\end{itemize}

The descriptions may be supported by manual annotations and optional perception-derived evidence.

\subsection{Layer 2: symbolic fact construction}

For the connector placement step $s_4$, Layer 2 may produce the following predicates with their respective confidence values:

\[
hasAction(s_4,\text{place}) \quad (0.90)
\]

\[
usesObject(s_4,o_5) \quad (0.88), \qquad usesObject(s_4,o_6) \quad (0.86)
\]

\[
isA(o_5,\text{Connector}) \quad (0.82), \qquad isA(o_6,\text{Terminal}) \quad (0.78)
\]

If geometric evidence is available, Layer 2 may also produce spatial predicates:

\[
aligned(o_5,o_6) \quad (0.75), \qquad touching(o_5,o_6) \quad (0.80).
\]

For the securing step $s_5$, the following predicates may be produced:

\[
hasAction(s_5,\text{secure}) \quad (0.87)
\]

\[
usesObject(s_5,o_5) \quad (0.85), \qquad usesObject(s_5,o_7) \quad (0.83)
\]

\[
isA(o_5,\text{Connector}) \quad (0.82), \qquad isA(o_7,\text{FasteningScrew}) \quad (0.79).
\]

These predicates form the local fact sets $\mathcal{P}_{s_4}$ and $\mathcal{P}_{s_5}$.

\subsection{Layer 3: rule-based inference}

For $s_4$, Layer 3 may apply a connector placement rule:

\[
\begin{aligned}
&hasAction(s_4,\text{place}) \wedge isA(o_5,\text{Connector}) \\
&\wedge isA(o_6,\text{Terminal}) \wedge aligned(o_5,o_6) \\
&\wedge touching(o_5,o_6) \\
&\Rightarrow produces(s_4, placed(o_5,o_6)).
\end{aligned}
\]

Using minimum aggregation, the confidence of the inferred effect is:

\[
\text{conf}(placed(o_5,o_6)) =
\min(0.90,0.82,0.78,0.75,0.80) = 0.75.
\]

Layer 3 may also infer an implicit assembly condition stating that contact must be maintained between the connector and the terminal:

\[
\begin{aligned}
&hasAction(s_4,\text{place}) \wedge isA(o_5,\text{Connector}) \\
&\wedge isA(o_6,\text{Terminal}) \wedge touching(o_5,o_6) \\
&\Rightarrow requires(s_4, contactMaintained(o_5,o_6)).
\end{aligned}
\]


The confidence of this inferred condition is:

\[
\text{conf}(contactMaintained(o_5,o_6)) =
\min(0.90,0.82,0.78,0.80) = 0.78.
\]

For $s_5$, Layer 3 may infer that securing the connector requires the connector to have already been placed on the terminal:

\[
\begin{aligned}
&hasAction(s_5,\text{secure}) \wedge isA(o_5,\text{Connector}) \\
&\Rightarrow requires(s_5, placed(o_5,o_6)).
\end{aligned}
\]


The confidence of this inferred precondition is:

\[
\text{conf}(requires(s_5, placed(o_5,o_6))) =
\min(0.87,0.82) = 0.82.
\]

Layer 3 may also infer that a fastening screw is required for securing:

\[
\begin{aligned}
&hasAction(s_5,\text{secure}) \wedge isA(o_5,\text{Connector}) \\
&\Rightarrow requiresTool(s_5,\text{FasteningScrew}).
\end{aligned}
\]



The confidence of this inferred tool requirement is:

\[
\text{conf}(requiresTool(s_5,\text{FasteningScrew})) =
\min(0.87,0.82) = 0.82.
\]

Since $s_5$ explicitly uses $o_7$ and $isA(o_7,\text{FasteningScrew})$ holds with confidence $0.79$, the required tool is supported by the available predicates.

\subsection{Layer 4: validation and enrichment}

Layer 4 validates and enriches both steps. For $s_4$, the placement action is supported by the action predicate, semantic predicates, and spatial predicates. The inferred effect $placed(o_5,o_6)$ has confidence $0.75$. If the acceptance threshold is $\tau_{acc}=0.70$, then the placement step can be classified as accepted.

For $s_5$, the inferred precondition $placed(o_5,o_6)$ can be linked to the effect produced by $s_4$. The required fastening screw is also supported by the predicates
\[
usesObject(s_5,o_7)
\]
and
\[isA(o_7,\text{FasteningScrew})
\]
The weakest relevant confidence is $0.79$, from the semantic classification of the fastening screw. If $\tau_{acc}=0.70$, then $s_5$ can also be classified as accepted.

If the spatial predicate $touching(o_5,o_6)$ were missing or had low confidence, $s_4$ could instead be classified as uncertain, because the placement effect would not be sufficiently supported. This illustrates how uncertainty in the evidence propagates into the validation decision without changing the rule structure.

\subsection{Graph representation and explanation trace}

The procedural graph preserves the input sequence:

\[
NEXT(s_4,s_5).
\]

A dependency edge is also added because $s_5$ requires the connector to be placed on the terminal, and this condition is produced by $s_4$:

\[
DEPENDS\_ON(s_5,s_4).
\]

The explanation trace for $s_5$ records the supporting evidence:

\begin{itemize}
    \item action evidence: $hasAction(s_5,\text{secure})$ with confidence $0.87$,
    \item object evidence: $usesObject(s_5,o_5)$ with confidence $0.85$,
    \item tool evidence: $usesObject(s_5,o_7)$ with confidence $0.83$,
    \item semantic evidence: $isA(o_5,\text{Connector})$ with confidence $0.82$ and $isA(o_7,\text{FasteningScrew})$ with confidence $0.79$,
    \item inferred precondition: $requires(s_5, placed(o_5,o_6))$ with confidence $0.82$,
    \item dependency support: $produces(s_4, placed(o_5,o_6))$ with confidence $0.75$.
\end{itemize}

This example shows that the same rule-based mechanism can be applied beyond the sleeve-and-rod example. The domain changes, but the reasoning pattern remains the same: step-level evidence is converted into predicates, rules infer missing procedural conditions, confidence values are propagated, and the resulting graph records both temporal order and dependency relations.






\chapter{Design}